%% Reduced-basis conservative residual inference for EXHALE steady winds.
%% Author: Kwang-Il Seon
\documentclass[english,a4paper]{my_memo}
\synctex=1
\usepackage{booktabs}
\usepackage{array}
\usepackage{amsmath}
\usepackage{babel}
\emergencystretch=3em

\newcommand{\code}[1]{\texttt{#1}}
\newcommand{\vect}[1]{\mathbf{#1}}
\newcommand{\Rres}{\vect{R}}
\newcommand{\Ustate}{\vect{U}}
\newcommand{\coef}{\vect{a}}
\newcommand{\flux}{\vect{F}}
\newcommand{\source}{\vect{S}}
\newcommand{\mdot}{\dot M}
\newcommand{\Rp}{R_{\rm p}}
\newcommand{\Mach}{\mathcal{M}}

\title{Reduced-Basis Steady Solver for EXHALE}
\author{Kwang-Il Seon}
\date{2026-08-23}

\begin{document}
\maketitle

\begin{quote}\small\itshape
This memo develops a direct steady-state solution strategy for EXHALE that
does not place the full, hours-long time-dependent calculation inside an
MCMC loop.  The proposed state is a smooth, low-dimensional expansion in
radius.  Mass conservation is imposed analytically, while momentum, energy,
boundary, chemistry, radiation, and sonic-point conditions are enforced by
scaled residuals.  Bulk advection, viscous and conductive transport, elemental
diffusion, and ion-stage advection are treated according to their conservative
fluxes and current status in EXHALE.  Two distinct formulations are documented:
(1) a reduced
basis evaluated by EXHALE's existing finite-volume residual, which targets
the current discrete fixed point exactly, and (2) a conservative weak
spectral-element formulation in which source terms are integrated by Gaussian
quadrature.  Nonlinear least squares is the primary solver.  Annealed MCMC,
sequential Monte Carlo, or posterior sampling is reserved for finding multiple
branches and measuring degeneracies.  This is a design document, not a claim
that the proposed solver is already implemented or validated.  The current
production workflow and the reduced proposal are compared explicitly.  Public
numerical precedents are also surveyed, with AUTO-07p selected as the primary
boundary-value reference and ESTER as the primary astrophysical reference.
\end{quote}

\tableofcontents
\bigskip

% =====================================================================
\section{Objective and scope}\label{sec:objective}
% =====================================================================

The computational problem is to obtain a physically correct stationary
solution of the one-dimensional radiation-hydrodynamic equations solved by
EXHALE without following the long transient required by explicit relaxation.
The intended use is a single planet and a fixed set of physical inputs, not a
Bayesian retrieval that reruns the complete hydrodynamic model at every sample.

Let the stationary equations be represented abstractly as
\begin{equation}
  \mathcal{R}[\Ustate;\boldsymbol{\theta}]=0,
  \label{eq:operator}
\end{equation}
where \(\Ustate(r)\) is the radial atmospheric state and
\(\boldsymbol{\theta}\) contains fixed planetary, stellar, atomic, and boundary
parameters.  A conventional mesh-based Newton solve treats all values of
\(\Ustate\) as independent unknowns.  The proposal instead restricts the state
to a smooth manifold,
\begin{equation}
  \Ustate(r) = \mathcal{P}(r;\coef), \qquad
  \coef\in\mathbb{R}^{K}, \qquad K\ll 3N,
  \label{eq:reducedmap}
\end{equation}
and searches for coefficients \(\coef\) that make the governing residual small.
The map \(\mathcal{P}\) must enforce positivity, composition normalization, and
as much exact conservation as possible before any optimizer is called.

The method is related to weighted-residual, collocation, spectral, and
spectral-element methods \cite{Finlayson1972,Canuto2006,Trefethen2000}.  A
residual interpreted as a statistical likelihood also resembles Bayesian
physics-informed and probabilistic differential-equation methods
\cite{Chkrebtii2016,Yang2021,Kramer2021}.  The spline implementation proposed
here is deliberately much smaller and more physically constrained than a
general neural network.

The design has four goals:
\begin{enumerate}
\item reduce the nonlinear search from values at every radial cell to a few
  tens of coefficients;
\item preserve the conservative structure of the hydrodynamic equations;
\item retain EXHALE's existing radiation, ionization, heating, and cooling
  calculations rather than replace them with a fitted surrogate; and
\item distinguish numerical solution finding from posterior uncertainty
  quantification, so that an arbitrary likelihood width is not mistaken for
  physical uncertainty.
\end{enumerate}

% =====================================================================
\section{The current EXHALE fixed point}\label{sec:current}
% =====================================================================

EXHALE inherits the Godunov finite-volume hydrodynamic core of ATES
\cite{Caldiroli2021}.  Its main marching path reconstructs left and right
states, evaluates numerical fluxes and sources, and applies a three-stage
strong-stability-preserving Runge--Kutta update.  The relevant implementation
is in
\begin{center}
\begin{tabular}{@{}p{0.38\textwidth}p{0.56\textwidth}@{}}
\toprule
file & role \\
\midrule
\code{src/EXHALE\_main.f90} & three-stage hydrodynamic update and the
  optional steady-solver hand-off \\
\code{src/modules/time\_step/RK\_rhs.f90} & spherical finite-volume flux
  divergence and hydrodynamic source terms \\
\code{src/modules/time\_step/}\newline
\code{steady\_residual.f90} & shared mass, momentum,
  and energy residual \\
\code{src/modules/time\_step/}\newline
\code{steady\_newton.f90} & state packing, equilibrium
  radiation/chemistry refresh, PTC, and JFNK \\
\code{src/modules/wind\_ae/wae\_difeq.f90} & a separate steady BVP relaxation
  with base and sonic-point conditions \\
\bottomrule
\end{tabular}
\end{center}

For a physical cell \(j\), the present residual has the schematic form
\begin{align}
 R_{\rho,j} &= (\nabla\!\cdot\!F_\rho)_j-S_{\rho,j},\\
 R_{p,j} &= (\nabla\!\cdot\!F_p)_j-S_{p,j}-S_{\mu,j},\\
 R_{E,j} &= (\nabla\!\cdot\!F_E)_j-S_{E,j}
             -(\mathcal{H}_j-\mathcal{C}_j)-S_{E,\mu,j},
 \label{eq:exhaleresidual}
\end{align}
where the optional \(\mu\) terms denote viscous work, molecular transport, and
thermal conduction.  The same reconstruction and flux calculation are used by
time marching and by the steady residual.  A true fixed point of the current
discretization therefore satisfies the same equation that the explicit code
relaxes.

Before evaluating Eq.~\eqref{eq:exhaleresidual}, the full residual path converts
the candidate conserved state to primitive variables, computes temperature,
updates excited hydrogen when enabled, solves the ionization equilibrium, and
recomputes heating and cooling.  This detail is essential: reducing only the
hydrodynamic unknowns does not freeze the microphysics.  It eliminates the
species fractions conditionally at every trial state.

The local Wind-AE implementation provides a second useful precedent.  It solves
a steady boundary-value problem directly, with density, temperature, velocity,
ionization, and column densities as mesh unknowns.  It imposes both
\(\Mach=1\) and the regularity numerator at the sonic point, then applies Newton
relaxation.  The reduced-basis proposal changes the representation and search
algorithm, but the base-to-sonic BVP structure is already present in the source
tree.  Murray-Clay et al.~\cite{MurrayClay2009} give the physical archetype for
this class of transonic escape solution.

% =====================================================================
\section{Comparison of the current and reduced solution strategies}
\label{sec:comparison}
% =====================================================================

The current EXHALE strategy and the proposed reduced strategy should not be
viewed as mutually exclusive solvers.  The production path is a conservative
and robust way to approach a discrete attractor, while the reduced path is a
candidate accelerator for reaching the neighborhood of a stationary root.  A
fair comparison must include the current hybrid workflow rather than describe
EXHALE as time marching alone.

\subsection{What the current workflow actually solves}

The production sequence is:
\begin{enumerate}
\item relax the finite-volume equations with three-stage SSP-RK updates;
\item when the mass-flux metric is sufficiently flat, optionally hand the full
  grid state to the JFNK steady solver;
\item evaluate equilibrium radiation, chemistry, heating, and cooling inside
  each nonlinear residual call;
\item include viscosity and thermal conduction through the same transport
  operator used by the marching path; and
\item when helium diffusion is active, alternate a JFNK hydrodynamic solve with
  500 implicit diffusion relaxation steps for as many as five outer passes.
\end{enumerate}
The last operation is an outer co-convergence, not a monolithic
hydrodynamics--diffusion Newton system.  The JFNK residual contains mass,
momentum, and energy but not the helium element-transport equation.  Ion-stage
advection is evaluated later by \code{post\_process\_adv.f90} and does not feed
back into the converged equilibrium hydrodynamics.

The reduced proposal instead maps a small coefficient vector to the radial
state, eliminates density by steady mass continuity where valid, refreshes the
same microphysics, and solves a stationary residual directly.  Depending on the
formulation, that residual is either the existing EXHALE finite-volume residual,
a strong collocation defect, or a conservative weak balance.

\begin{center}\small
\begin{tabular}{@{}p{0.21\textwidth}p{0.34\textwidth}p{0.37\textwidth}@{}}
\toprule
property & current EXHALE workflow & reduced-basis workflow \\
\midrule
unknowns & conserved state in every radial cell, approximately \(3N\)
  hydrodynamic unknowns & typically tens of spline coefficients plus global
  quantities such as \(\mdot\) and possibly a sonic radius \\
route to stationarity & explicit relaxation followed by an optional full-grid
  PTC/JFNK finish & direct nonlinear solve of the stationary coefficient
  residual, normally with continuation \\
conservation & face-based finite-volume conservation using the production
  HLLC/WENO reconstruction & exact only for constraints built into the map or
  for correctly assembled conservative element residuals \\
sharp structure & robust numerical-flux treatment of shocks and narrow layers &
  smooth and efficient when resolved, but susceptible to missed layers or
  oscillatory basis error \\
sonic branch & approached through characteristic time evolution and polished by
  JFNK & must be selected by explicit topology, regularity, and boundary
  constraints \\
species diffusion & implicit operator splitting and outer hydro--diffusion
  co-convergence & can retain the production outer loop or add element fractions
  and flux equations to a simultaneous residual \\
ion-stage advection & guarded post-processing without hydrodynamic feedback &
  requires additional transported fields or a conditional upstream solve \\
discrete target & unambiguously the production finite-volume fixed point & the
  same fixed point only in Formulation A; collocation and weak forms define
  different discrete problems \\
\bottomrule
\end{tabular}
\end{center}

\subsection{Advantages of the current EXHALE workflow}

\paragraph{Conservation and a defined fixed point.}
The same reconstruction, numerical face flux, boundary routine, and source
calculation are used in marching and in \code{steady\_residual.f90}.  A
converged JFNK state therefore solves the equation that the production
discretization relaxes.  This supplies an unambiguous acceptance diagnostic.

\paragraph{Robustness to localized and nonsmooth structure.}
The finite-volume path is designed for steep fronts, contact modes, and shocks.
It is safer than a coarse smooth basis when an ionization front, thermal layer,
or diffusion scale height is narrower than the proposed basis support.

\paragraph{A comparatively broad basin of attraction.}
Time evolution follows causal propagation and physical relaxation.  It can
bring an imperfect initial atmosphere into the local convergence basin of JFNK.
The code also returns to time marching if the JFNK finish fails, rather than
accepting an unconverged Newton state.

\paragraph{Reuse of established microphysics.}
Radiation, chemistry, heating, cooling, viscosity, conduction, reconstruction,
and boundaries remain on the normal production path.  There is no second
implementation whose agreement must be assumed.

\subsection{Disadvantages of the current EXHALE workflow}

\paragraph{Long transients.}
The calculation may spend most of its runtime evolving slow thermal, chemical,
or diffusive modes even when only one stationary state is required.  These
hours-long trajectories are computational work toward the root, not additional
information needed by a steady calculation.

\paragraph{A high-dimensional nonlinear search.}
The full-grid JFNK solve retains approximately \(3N\) hydrodynamic directions,
including many radial perturbations that a smooth wind does not need.  The
resulting Jacobian can be poorly conditioned, and GMRES performance depends on
the warm state and preconditioning.

\paragraph{Local convergence of the Newton finish.}
The production code deliberately waits until the mass-flux metric reaches the
handoff neighborhood.  JFNK accelerates the final convergence but does not
currently replace the global relaxation stage.

\paragraph{Split composition transport.}
Helium diffusion is co-converged outside the JFNK residual, while ion-stage
advection remains post-processing.  This design is practical and internally
documented, but it does not provide one Jacobian or one simultaneous residual
for all transported composition fields.

\paragraph{Limited branch tracing.}
One time evolution normally selects one attractor.  Parameter folds and several
stationary branches are harder to map systematically than with
pseudo-arclength continuation.

\subsection{Advantages of the reduced strategy}

\paragraph{Removal of redundant search directions.}
A smooth wind can often be represented by tens of coefficients rather than
values in every cell.  Exact mass-continuity elimination removes density as an
independent function and gives \(\mdot\) a direct global meaning.

\paragraph{Direct stationary solve.}
The method seeks \(\vect{r}(\coef)=0\) without integrating the complete transient.
It can therefore avoid the slowest relaxation modes.  This is a potential
speed advantage, not yet a measured EXHALE result: every trial may still require
expensive nonlocal radiation and equilibrium chemistry.

\paragraph{Constraints built into the representation.}
Positive velocity and temperature, normalized composition, fixed base values,
and exact steady mass continuity can be imposed before the nonlinear solver is
called.  This prevents exploration of many invalid full-grid perturbations.

\paragraph{Natural continuation.}
Radiation, chemistry, conduction, and diffusion can be activated progressively.
AUTO-style pseudo-arclength continuation can also trace folds and distinguish
subsonic, breeze, and transonic branches more systematically than a single
relaxation trajectory.

\paragraph{Potential simultaneous transport solution.}
Adding a small abundance basis and a conservative total element-flux residual
can place hydrodynamics and diffusion in one nonlinear system.  This is useful
only after the target energy and momentum closures have been stated; adding
diffusive enthalpy or frictional heating would define physics beyond the current
operator-split model.

\subsection{Disadvantages and risks of the reduced strategy}

\paragraph{Basis truncation is a new model error.}
A small scalar merit does not prove that a narrow heating, cooling, ionization,
or diffusion layer is represented.  Basis refinement, an independently located
held-out residual, and comparison with the production-grid residual are all
required.

\paragraph{Conservation is not automatic.}
A small strong residual at collocation nodes does not by itself guarantee an
accurate integrated flux balance.  Advective and diffusive contributions must
enter one total conservative species flux.  Weak conservation improves the
integrated balance but no longer targets the exact HLLC/WENO fixed point.

\paragraph{Sonic regularity becomes an explicit responsibility.}
The reduced BVP can converge to a smooth subsonic breeze unless the intended
branch and critical conditions are imposed.  A differential equation divided by
\(v^2-c_s^2\) is singular at the sonic point and must be desingularized, split
into domains, or accompanied by an explicit regularity equation.

\paragraph{A small state does not imply a cheap residual.}
Spline coefficients influence several cells, and column-dependent radiation
couples the radial domain nonlocally.  The reduced Jacobian is smaller but may be
dense, while one residual evaluation may retain nearly the full chemistry cost.

\paragraph{Smooth bases are unsuitable for genuine discontinuities.}
A shock requires domain decomposition, an explicit shock location, or retention
of the finite-volume solver.  Regularization must not smooth away a physically
required discontinuity merely to improve the scalar objective.

\subsection{Recommended hybrid architecture}

The safest use of dimension reduction is as a front end to the production
solver:
\begin{equation}
 \begin{gathered}
  \text{simple Parker or frozen-source state}
  \longrightarrow \text{reduced continuation}
  \longrightarrow \text{reduced Formulation A solve} \\
  \longrightarrow \text{full-grid JFNK correction}
  \longrightarrow \text{short production evolution and residual check}.
 \end{gathered}
 \label{eq:hybridworkflow}
\end{equation}
The first reduced prototype should use the existing finite-volume residual, not
a newly derived continuum residual.  It should solve only for \(v\), \(T\), and
\(\mdot\), with density eliminated analytically and equilibrium chemistry
refreshed conditionally.  After that path is verified, adaptive collocation,
weak conservation, and simultaneous elemental diffusion can be introduced in
separate validation stages.

The current solver remains the final physical diagnostic and robust fallback.
The reduced solver is successful if it places the state inside the production
JFNK basin substantially faster while preserving the final EXHALE residual and
physical branch.  Wall-clock acceleration must be measured; it must not be
inferred solely from the smaller coefficient count.

% =====================================================================
\section{Two formulations that must not be conflated}\label{sec:twoforms}
% =====================================================================

There are two technically different ways to combine a reduced basis with
EXHALE.  They have different targets and different validation requirements.

\subsection{Formulation A: reduced state, existing discrete residual}

The coefficient vector is mapped onto all \(N\) EXHALE cells,
\begin{equation}
  \coef \xrightarrow{\mathcal{P}_N} \Ustate_N(\coef)
  \xrightarrow{\text{existing EXHALE evaluator}}
  \vect{F}_N(\coef)\in\mathbb{R}^{3N}.
  \label{eq:formA}
\end{equation}
The reduced solve minimizes a scaled norm of \(\vect{F}_N\).  This formulation
still evaluates every cell, but one evaluation is a stationary residual call,
not a full time integration.  It has three decisive advantages:
\begin{enumerate}
\item it targets the exact finite-volume fixed point of the current code;
\item all existing boundary, WENO3/HLLC, radiation, chemistry, and transport
  paths are reused; and
\item the full-cell residual exposes localized errors that a quadrature grid
  could miss.
\end{enumerate}

It is the safest prototype because a failure cannot be attributed to a newly
derived continuum residual.  Gaussian quadrature is not needed in this first
formulation.

\subsection{Formulation B: conservative weak residual with quadrature}

The second formulation defines a new spatial discretization.  Divide the radial
domain into \(N_e\) elements \(I_e=[r_e,r_{e+1}]\).  If
\(A(r)=r^2\), omitting the common solid-angle factor, a stationary conservation
law reads
\begin{equation}
  \frac{d}{dr}\left[A(r)\flux(\Ustate)\right]
  = A(r)\source(\Ustate,r).
  \label{eq:conservative}
\end{equation}
Integrating over one element gives the conservative weak residual
\begin{equation}
  \Rres_e = A_{e+1}\flux(\Ustate_{e+1})
            -A_e\flux(\Ustate_e)
            -\int_{r_e}^{r_{e+1}}A(r)\source(\Ustate,r)\,dr.
  \label{eq:weakresidual}
\end{equation}
With \(Q\)-point Gauss--Legendre quadrature,
\begin{equation}
 \int_{r_e}^{r_{e+1}}A\source\,dr
 \simeq \frac{\Delta r_e}{2}
 \sum_{q=1}^{Q}w_q A(r_{e,q})
 \source\!\left[\mathcal{P}(r_{e,q};\coef),r_{e,q}\right].
 \label{eq:gauss}
\end{equation}

Equation~\eqref{eq:weakresidual} conserves the integrated flux exactly up to
quadrature and nonlinear-solve errors.  It does not require noisy numerical
derivatives of the source terms.  It is preferable to minimizing the squared
pointwise differential residual when the atmosphere contains narrow heating,
cooling, or ionization transitions.  However, it no longer solves the exact
HLLC/WENO3 finite-volume equation.  Agreement with Formulation A should be
expected only after basis, element, and quadrature convergence.

\subsection{Recommendation}

Implement Formulation A first.  It tests whether dimension reduction itself
removes the harmful search directions.  Implement Formulation B only after the
reduced manifold is demonstrated to contain the desired discrete solution.
The two results should be reported as a discrete EXHALE fixed point and a weak
continuum solution, respectively.

% =====================================================================
\section{Low-dimensional representation}\label{sec:basis}
% =====================================================================

\subsection{Radial coordinate}

Map the physical domain to \(x\in[0,1]\).  A logarithmic map is usually more
appropriate than a linear map because most thermal and ionization structure is
near the lower boundary:
\begin{equation}
 x(r)=\frac{\ln(r/r_{\min})}{\ln(r_{\max}/r_{\min})}.
 \label{eq:coord}
\end{equation}
For an L1-truncated Roche domain, the map remains finite and monotone, but the
outer boundary must be treated according to its actual characteristic structure;
it must not be called a supersonic outflow boundary when the solution is
subsonic there.

\subsection{Velocity, temperature, and mass flux}

For an outward solution with \(v>0\), use cubic B-splines
\begin{align}
 \ln v(x) &= \sum_{k=1}^{K_v}a_k B_k(x),\\
 \ln T(x) &= \sum_{k=1}^{K_T}b_k B_k(x),\\
 m &= \ln\mdot.
 \label{eq:splines}
\end{align}
The logarithms enforce positive velocity, temperature, and mass-loss rate.  A
single global Chebyshev expansion is attractive for very smooth Parker winds,
but local B-splines are safer for EXHALE because their support is limited and
knots can be concentrated near the base, a heating peak, an ionization front,
or the sonic point.  Global polynomials may oscillate across a narrow transition.

If a solution may contain inflow or a stagnation point, \(\ln v\) is invalid.
One alternative is
\begin{equation}
  v(x)=v_{\rm sc}\sinh\!\left[\sum_k a_kB_k(x)\right],
  \label{eq:signedv}
\end{equation}
which permits a sign change without a hard kink.  The mass-flux elimination
below is then singular at \(v=0\), so a steady solution with nonzero \(\mdot\)
cannot cross a true stagnation point.  Such a state belongs to a different
physical branch and should not be hidden by regularization.

\subsection{Exact elimination of steady mass continuity}

For a steady spherical single-fluid wind,
\begin{equation}
  4\pi r^2\rho(r)v(r)=\mdot.
  \label{eq:masscontinuity}
\end{equation}
Therefore
\begin{equation}
  \rho(r;\coef)=\frac{\exp(m)}{4\pi r^2v(r;\coef)}.
  \label{eq:rhoeliminate}
\end{equation}
Density is not assigned an independent spline.  Equation
\eqref{eq:rhoeliminate} removes an entire function, prevents a flat but
radially inconsistent mass flux, and gives \(\mdot\) a direct physical meaning.
The factor \(4\pi\) can be omitted internally if EXHALE's code-unit flux is used,
provided the convention is applied consistently.

This elimination assumes the mass equation has no distributed source.  Any
future mass loading, condensation, or separate-fluid transfer would invalidate
Eq.~\eqref{eq:masscontinuity}; in that case the mass equation must return to the
residual system.

\subsection{Pressure and composition}

Given \(\rho\), \(T\), and equilibrium species fractions, pressure follows from
the same equation of state used by EXHALE,
\begin{equation}
  p = (n_{\rm tot}+n_e)k_{\rm B}T.
  \label{eq:eos}
\end{equation}
The coefficient-to-state map must call the existing composition and equation-of-
state routines.  A constant mean molecular weight would change the physical
problem and would be acceptable only for a hydrodynamics-only prototype that is
explicitly labeled as such.

The preferred design eliminates ion fractions through the existing equilibrium
solver at every trial state.  If an advected species must instead be represented,
use logits.  For stages \(s=1,\ldots,S\),
\begin{equation}
 y_s(x)=\frac{\exp z_s(x)}{\sum_{t=1}^{S}\exp z_t(x)},
 \qquad z_s(x)=\sum_k c_{sk}B_k(x),
 \label{eq:softmax}
\end{equation}
which enforces \(0<y_s<1\) and \(\sum_sy_s=1\).  One logit should be fixed to
zero to remove the redundant common offset.

\subsection{Initial coefficient count}

A practical first trial is
\begin{center}
\begin{tabular}{lcc}
\toprule
quantity & coefficients & treatment \\
\midrule
\(v(r)\) & 8--12 & cubic B-spline in \(\ln v\) \\
\(T(r)\) & 8--12 & cubic B-spline in \(\ln T\) \\
\(\mdot\) & 1 & scalar in \(\ln\mdot\) \\
\(r_s\) & 0 or 1 & explicit only when a sonic constraint is used \\
\(\rho(r)\) & 0 & eliminated by Eq.~\eqref{eq:rhoeliminate} \\
equilibrium species & 0 & solved conditionally at residual nodes \\
\bottomrule
\end{tabular}
\end{center}
Thus the first nonlinear search has approximately 17--26 variables.  Start with
the coarsest basis that can represent a hot lower atmosphere and one sonic
transition, then add knots by continuation.  A small residual obtained only
after adding many nearly redundant coefficients is evidence of poor scaling or
an unresolved physical layer, not evidence that a larger MCMC is required.

% =====================================================================
\section{Residual objective}\label{sec:objectivefun}
% =====================================================================

\subsection{Scaled physics merit}

Let \(\vect{r}(\coef)\) collect hydrodynamic, boundary, sonic, and optional
chemistry residuals.  Define
\begin{equation}
 \Phi(\coef)=\frac{1}{2}\left\|\vect{W}\vect{r}(\coef)\right\|_2^2
             +\Phi_{\rm prior}(\coef).
 \label{eq:merit}
\end{equation}
The diagonal or block matrix \(\vect{W}\) nondimensionalizes each residual.
Without physical scaling, base momentum can dominate merely because gravity and
pressure individually contain large cancelling terms, while an energy imbalance
elsewhere is numerically smaller in code units.

For Formulation A, retain and extend EXHALE's cell-local state scales.  A useful
starting point is
\begin{align}
 s_{\rho,j} &= \max(\rho_j,\rho_{\rm floor})/t_{\rm dyn,j},\\
 s_{p,j} &= \rho_j(|v_j|+c_{s,j})/t_{\rm dyn,j},\\
 s_{E,j} &= \max(E_j,p_j,E_{\rm floor})/t_{\rm dyn,j},
 \label{eq:scales}
\end{align}
with \(t_{\rm dyn,j}\) tied to the local crossing time.  The exact scales must
be tested by inspecting contributions from every equation and radius.  No one
row should become invisible because another row contains a large cancellation.

For Formulation B, scale an element residual by a characteristic flux through
that element, not by its width alone:
\begin{equation}
 \widehat R_{k,e}=\frac{R_{k,e}}
 {\max(|A_eF_{k,e}|,|A_{e+1}F_{k,e+1}|,S_{k,e}^{\rm ref})}.
 \label{eq:weakscale}
\end{equation}
Floors must represent a physical hydrostatic or thermal scale.  A floor imported
from an unrelated cell can distort the solution.

\subsection{Boundary conditions}

Boundary conditions may be imposed in either of two ways.  A hard construction
builds them into \(\mathcal{P}\), for example
\begin{equation}
 \ln T(x)=\ln T_0+x\sum_k b_k\widetilde B_k(x),
 \label{eq:hardbc}
\end{equation}
which guarantees \(T(0)=T_0\).  A soft condition adds
\begin{equation}
 \Phi_{\rm BC}=\frac{1}{2}\sum_i
 \left[\frac{B_i(\Ustate)}{s_{B_i}}\right]^2.
 \label{eq:softbc}
\end{equation}
Hard conditions reduce dimension and avoid weight tuning, but they should be
used only for genuinely fixed physical conditions.  A hard base velocity would
conflict with EXHALE's valve or mass-flux condition.  The base density and
temperature, pressure-gradient consistency, and outer characteristic condition
must be copied from the selected EXHALE configuration rather than inferred from
an obsolete memo.

\subsection{Sonic-point conditions}

The conservative equations are not algebraically singular at \(\Mach=1\), but
the transonic branch is selected by regularity.  If the sonic radius \(r_s\) is
an unknown, add
\begin{align}
 R_{\Mach} &= \Mach(r_s)-1,\\
 R_{\rm crit} &= N_{\rm wind}(r_s),
 \label{eq:sonicconditions}
\end{align}
where \(N_{\rm wind}\) is the numerator of the differential wind equation after
the sonic denominator has been separated.  The sonic merit is
\begin{equation}
 \Phi_s=\frac{1}{2}\left(\frac{R_{\Mach}}{s_{\Mach}}\right)^2
       +\frac{1}{2}\left(\frac{R_{\rm crit}}{s_{\rm crit}}\right)^2.
 \label{eq:sonicmerit}
\end{equation}
This mirrors the Mach-one and numerator conditions in
\code{wae\_difeq.f90}.  For a domain whose physical solution is subsonic at L1,
forcing an interior sonic point is wrong.  Branch topology must be chosen before
the constraint is enabled.

\subsection{Smoothness prior}

Spline truncation already imposes smoothness.  If additional regularization is
needed, use curvature rather than coefficient magnitude:
\begin{equation}
 \Phi_{\rm smooth}=\frac{\lambda_v}{2}
 \int_0^1\left[\frac{d^2\ln v}{dx^2}\right]^2dx
 +\frac{\lambda_T}{2}
 \int_0^1\left[\frac{d^2\ln T}{dx^2}\right]^2dx.
 \label{eq:smoothness}
\end{equation}
The same Gaussian quadrature can evaluate these integrals.  Regularization may
stabilize an underdetermined coarse solve, but it must not erase a physically
required ionization or cooling front.  A residual decrease obtained by
increasing \(\lambda\) is not convergence of the governing equations.

\subsection{Held-out residual}

Never assess the weak solution only at the quadrature nodes used in
Eq.~\eqref{eq:gauss}.  Evaluate a held-out residual on a denser, independently
located grid.  A low training residual and a large held-out residual indicate
aliasing or a basis that oscillates between nodes.  For Formulation B, also
double \(Q\) without refitting and require the integrated balances to remain
stable.

% =====================================================================
\section{Radiation and chemistry}\label{sec:microphysics}
% =====================================================================

\subsection{Conditional elimination}

At a trial \(v(r)\), \(T(r)\), and \(\mdot\), the density is known from
Eq.~\eqref{eq:rhoeliminate}.  The preferred residual evaluation is then:
\begin{enumerate}
\item construct \(\rho\), \(v\), and \(T\) at all required nodes;
\item compute cumulative columns and attenuated radiation consistently;
\item solve the coupled equilibrium species system at those nodes;
\item recompute \(n_{\rm tot}+n_e\), pressure, excited-state populations,
  heating, and cooling; and
\item assemble hydrodynamic residuals from this self-consistent state.
\end{enumerate}
This is a nonlinear variable-elimination or variable-projection strategy.  It
keeps chemical unknowns out of \(\coef\) while retaining their response to the
candidate hydrodynamic state.

\subsection{Nonlocal columns}

Photoionization couples radii through outward columns.  Formulation A should use
the current EXHALE column calculation unchanged.  Formulation B has two choices:
\begin{enumerate}
\item compute each column by cumulative Gaussian quadrature from the outer
  boundary; or
\item introduce column fields satisfying
  \begin{equation}
   \frac{dN_s}{dr}=-n_s(r), \qquad N_s(r_{\max})=N_{s,\max},
   \label{eq:columnode}
  \end{equation}
  and add their weak residuals.
\end{enumerate}
The first choice has fewer unknowns and is the recommended start.  The second
choice localizes the radiation equations and may provide a sparser Jacobian,
but it adds at least one field for every opacity carrier.

\subsection{Determinism requirement}

The present full residual updates the species fraction array in place.  This is
efficient inside Newton continuation, but a statistical likelihood must be a
deterministic function of \(\coef\).  If the same \(\coef\) yields different
residuals depending on the previous proposal's chemistry, detailed balance is
violated and even deterministic optimization can acquire a false noise floor.

Each trial must therefore use one of the following policies:
\begin{enumerate}
\item restore the same reference species state before every evaluation and
  converge equilibrium tightly;
\item prove by repeated evaluations that the local equilibrium is unique and
  independent of its initial guess; or
\item include the relevant chemical branch in the state vector explicitly.
\end{enumerate}
Cache keys must include every state variable on which a rate depends.  Caching
based only on radius or column would make coefficient perturbations return stale
physics.

\subsection{Differentiability}

Hard switches, clipped fractions, \(\max(v,0)\) valves, table boundaries, and
fallback chemistry solvers make \(\vect{r}(\coef)\) nonsmooth.  This matters
more for Gauss--Newton or Hamiltonian Monte Carlo than for time marching.  The
reduced solver should use the existing smooth valve where physically permitted,
record every branch switch, and reject proposals that trigger an unconverged
chemistry result.  Smoothing a switch is acceptable only when its validity range
and effect on the physical solution are documented.

% =====================================================================
\section{Advection and diffusion in the reduced formulation}
\label{sec:transport}
% =====================================================================

The reduced method remains applicable when advection and diffusion are active.
They do not require a return to one unknown at every cell, but they change which
fields, boundary conditions, and conservation equations must be retained.  The
essential rule is to reduce the state representation, not the conservative
transport law.

\subsection{Transport already present and transport still outside the fixed point}

The present EXHALE paths must be distinguished before defining the target
residual:
\begin{center}\small
\begin{tabular}{@{}p{0.25\textwidth}p{0.28\textwidth}p{0.39\textwidth}@{}}
\toprule
process & present treatment & implication for a reduced solve \\
\midrule
bulk hydrodynamic advection & conservative finite-volume mass, momentum, and
  energy fluxes & already included in Formulation A; retain conservative fluxes
  in Formulation B \\
viscosity and thermal conduction & common operator used by marching and the
  steady residual when enabled & already included in the current discrete fixed
  point \\
elemental diffusion & implicit operator-split update, followed by outer
  hydrodynamics--diffusion co-convergence in the current Newton workflow & not
  a row of the JFNK residual; a direct simultaneous solve needs an additional
  conservative element-transport residual \\
ion-stage advection & guarded post-processing calculation & does not feed back
  into the current equilibrium hydrodynamic solution; simultaneous coupling
  defines an enlarged fixed point \\
\bottomrule
\end{tabular}
\end{center}

Thus, Formulation A already includes bulk advection, viscosity, and conduction
through the production residual.  It reproduces the fully co-converged result
with elemental diffusion only if the same outer iteration is retained.  It does
not become a simultaneous diffusion or ion-advection solve merely by evaluating
\code{steady\_residual.f90}.

\subsection{Bulk advection, viscosity, and conduction}

Bulk advection is not a new free function.  It is contained in the conservative
hydrodynamic flux \(\flux(\Ustate)\), while steady mass continuity still permits
the exact elimination in Eq.~\eqref{eq:rhoeliminate}.  In a continuum weak
formulation, viscosity and thermal conduction enter through the total momentum
and energy fluxes.  Schematically,
\begin{align}
 F_{p}^{\rm tot} &= F_{p}^{\rm adv}+p-F_{p}^{\rm visc}, \\
 F_{E}^{\rm tot} &= F_{E}^{\rm adv}+F_{E}^{\rm cond}
                    +F_{E}^{\rm visc},
 \label{eq:totalhydroflux}
\end{align}
with signs fixed by the production convention.  The weak residual uses these
total fluxes in the boundary term of Eq.~\eqref{eq:weakresidual}; the flux
derivatives should not be reintroduced as unrelated pointwise sources.  Cubic
B-splines provide the first derivatives needed by standard constitutive laws
and are normally sufficient.  Sharp conductive fronts still require local knot
or element refinement.

\subsection{Elemental diffusion as a conservative flux constraint}

When helium or a metal element separates from hydrogen, a constant elemental
abundance is no longer a valid conditional variable.  Introduce a bounded,
low-dimensional elemental fraction, for example
\begin{equation}
 \operatorname{logit}\!\left(\frac{f_{\rm He}}{f_{\max}}\right)
 =\sum_{k=1}^{K_{\rm He}}c_k B_k(x).
 \label{eq:heparameterization}
\end{equation}
Six to ten coefficients are a reasonable initial experiment, but the accepted
number must be determined by residual and basis refinement rather than fixed in
advance.

The helium transport implemented by EXHALE can be represented schematically by
one total flux,
\begin{equation}
 J_{\rm He}=n_{\rm He}v
 -n_{\rm H}D_{\rm He}
 \left[\frac{df_{\rm He}}{dr}
 +f_{\rm He}\left(G_{\rm He}
 +\alpha_T\frac{d\ln T}{dr}\right)\right]
 -n_{\rm H}K_{zz}\frac{df_{\rm He}}{dr},
 \label{eq:heflux}
\end{equation}
where \(G_{\rm He}\) collects the gravitational and ambipolar separation
terms, \(D_{\rm He}\) is the molecular diffusion coefficient, and \(K_{zz}\)
is the eddy diffusion coefficient.  The exact code expression, including its
effective mass and sign conventions, must be called or factored from
\code{species\_diffusion.f90}; Eq.~\eqref{eq:heflux} is not a replacement
implementation.  Multispecies escape and helium fractionation models provide
the broader physical context for coupled advective--diffusive fluxes
\cite{Koskinen2013a,Xing2023,Taylor2025}.

If chemical reactions conserve helium nuclei, steady elemental conservation is
\begin{equation}
 \frac{1}{r^2}\frac{d}{dr}\left(r^2J_{\rm He}\right)=0.
 \label{eq:heconservation}
\end{equation}
For an element \(e=[r_L,r_R]\) and a test function \(w_j\), its weak residual is
\begin{equation}
 R_{{\rm He},e,j}=
 \left[w_j r^2J_{\rm He}\right]_{r_L}^{r_R}
 -\int_{r_L}^{r_R}\frac{dw_j}{dr}\,r^2J_{\rm He}\,dr.
 \label{eq:heweak}
\end{equation}
Gaussian quadrature evaluates the remaining integral.  With a constant test
function, Eq.~\eqref{eq:heweak} is simply the difference between the total
element fluxes at the two faces.  This is the most transparent conservation
diagnostic.

Diffusion also requires boundary conditions that pure hydrodynamics does not.
The present physical choice is a fixed elemental ratio at the base and zero
\emph{diffusive} flux at the outer boundary.  The latter must not set the total
flux to zero: the outward advective flux remains allowed.  Applying separate
residuals to advective and diffusive divergences can create a numerical source
or sink; they must be combined into \(J_{\rm He}\) before conservation is
tested.

\subsection{Ion-stage advection}

Local ionization equilibrium can be eliminated conditionally only when reaction
times are short compared with transport times.  If ion stages are advected
self-consistently, each retained stage satisfies
\begin{equation}
 \frac{1}{r^2}\frac{d}{dr}\left(r^2J_{X^i}\right)
 =P_{X^i}-L_{X^i},
 \label{eq:ionadvection}
\end{equation}
where \(J_{X^i}\) includes the selected advective and diffusive contributions.
Its weak residual is
\begin{align}
 R_{X^i,e,j}={}&
 \left[w_jr^2J_{X^i}\right]_{r_L}^{r_R}
 -\int_{r_L}^{r_R}\frac{dw_j}{dr}\,r^2J_{X^i}\,dr \nonumber\\
 &-\int_{r_L}^{r_R}w_jr^2(P_{X^i}-L_{X^i})\,dr.
 \label{eq:ionweak}
\end{align}
The production and loss integrals are evaluated at Gaussian nodes using the
same attenuated radiation and thermochemical state as the hydrodynamic source
terms.

Two reductions are practical:
\begin{enumerate}
\item Represent \(S-1\) independent stage logits with splines and obtain all
  \(S\) fractions through a softmax map.  This supports a coupled boundary-value
  solve and preserves positivity and normalization, but adds coefficients.
\item For a trial hydrodynamic and elemental state, integrate the ion-advection
  equations conditionally from the upstream boundary.  This uses fewer global
  coefficients and is attractive for a strictly outward flow with \(v>0\).
\end{enumerate}
The second option is not generally valid through stagnation or inflow, because
the upstream boundary changes.  Such cases require a boundary-value treatment
or a conservative flux residual with characteristic boundary conditions.

Do not impose all ion-stage equations together with an independent equation for
their sum.  For one element, use one elemental transport equation and only
\(S-1\) independent stage equations.  When reaction bookkeeping exactly
conserves nuclei, the sum of the stage source terms is zero and the summed
stage residual must reproduce the elemental residual.  This identity is both a
rank test and a stringent implementation diagnostic.

\subsection{Energy consistency and the definition of the target problem}

A fully coupled multicomponent model may require more than composition fluxes.
Species diffusion can carry enthalpy,
\begin{equation}
 F_E^{\rm species}=\sum_s h_sJ_s^{\rm diff},
 \label{eq:speciesenthalpy}
\end{equation}
and interspecies friction can heat the gas.  Omitting these terms while adding
composition diffusion is an approximation whose validity range must be stated.
If the goal is the current EXHALE result, retain the present operator-split
physics and reproduce its outer co-convergence.  Adding
Eq.~\eqref{eq:speciesenthalpy} or new friction terms defines a different
physical model and requires independent conservation tests.

For a simultaneous reduced transport solve, the residual therefore has the
block structure
\begin{equation}
 \vect{r}_{\rm total}=
 \begin{bmatrix}
  \vect{r}_{\rm mass} & \vect{r}_{\rm momentum} & \vect{r}_{\rm energy} &
  \vect{r}_{\rm element} & \vect{r}_{\rm ion}
 \end{bmatrix}^{\mathsf T}.
 \label{eq:transportblocks}
\end{equation}
Blocks are included only for fields that are solved dynamically.  Equilibrium
ion fractions must not also receive advection residuals.  The extra spline
coefficients remain compatible with optimization and MCMC because their number
scales with the number of transported fields, not with the number of radial
cells.

% =====================================================================
\section{Optimization and Bayesian interpretation}\label{sec:inference}
% =====================================================================

\subsection{Why MCMC is not the primary root finder}

If the equations and boundary conditions define one deterministic solution,
the target is \(\vect{r}(\coef)=0\).  Nonlinear least squares uses residual
magnitudes and derivatives to approach that root.  A random-walk MCMC sampler
continues to explore after it has found the basin and therefore spends most
evaluations estimating a distribution that was introduced artificially.

The recommended hierarchy is:
\begin{enumerate}
\item obtain a coarse solution by trust-region least squares or
  Levenberg--Marquardt;
\item continue in basis size, radiation coupling, or source strength;
\item use multistart, differential evolution, or annealed sampling only when
  several basins are plausible; and
\item sample a posterior only after a stable minimum and a defensible error
  model have been established.
\end{enumerate}

\subsection{Pseudo-posterior for branch exploration}

A residual-based distribution can be defined as
\begin{equation}
 \pi_\beta(\coef)\propto
 \exp[-\beta\Phi(\coef)]\,p_0(\coef),
 \label{eq:pseudoposterior}
\end{equation}
where \(p_0\) supplies bounds and smoothness information.  Increasing
\(\beta\) concentrates the distribution around minima.  Parallel tempering or
sequential Monte Carlo can move between branches more effectively than one cold
chain.  The final modes provide starting points for deterministic root solves.

This construction is a search distribution, not automatically a physical
posterior.  Its width depends on \(\beta\), residual scaling, element count,
and regularization.  Bayesian physics-informed models make a similar residual-
likelihood construction \cite{Yang2021}, while probabilistic ODE/BVP methods
place probability on solution functions and discretization error
\cite{Chkrebtii2016,Kramer2021}.  Those frameworks justify uncertainty only
after the stochastic model has been specified; the exponential form alone does
not do so.

\subsection{A physical posterior}

A defensible posterior requires uncertain physical inputs, model discrepancy,
or observational data.  One possible hierarchy is
\begin{equation}
 p(\coef,\boldsymbol{\theta}\mid D)\propto
 p(D\mid\mathcal{G}[\mathcal{P}(\coef),\boldsymbol{\theta}])
 p(\vect{r}(\coef,\boldsymbol{\theta})\mid\Sigma_{\rm model})
 p_0(\coef,\boldsymbol{\theta}),
 \label{eq:physicalposterior}
\end{equation}
where \(D\) is an observed spectrum, \(\mathcal{G}\) is the transit model, and
\(\Sigma_{\rm model}\) has a stated interpretation.  Hydrodynamic escape
analyses have used MCMC with grids or fast Parker-wind models to infer
temperature, composition, and mass loss \cite{Lampon2021,DosSantos2022}.
Those are inference examples, not examples of replacing a hydrodynamic root
solver with MCMC.

\subsection{Sampler choice}

For \(K\lesssim30\), an affine-invariant ensemble sampler is easy to deploy and
parallelize \cite{Goodman2010}, but it does not solve multimodality by itself.
Hamiltonian Monte Carlo can be much more efficient in a smooth correlated basin,
but requires reliable gradients.  Function-space MCMC methods become relevant
only if basis refinement is taken toward a high-dimensional random function
\cite{Cotter2013}.  For the proposed small spline expansion, careful scaling
and a deterministic optimizer are more important than a specialized
infinite-dimensional sampler.

Convergence reporting must include more than acceptance fraction.  At minimum
report independent-chain agreement, integrated autocorrelation times, effective
sample sizes, trace plots for \(\ln\mdot\) and residual components, mode
occupancy, and repeated likelihood determinism tests.

% =====================================================================
\section{Jacobians and computational cost}\label{sec:jacobian}
% =====================================================================

Let \(\vect{J}_a=\partial\vect{r}/\partial\coef\) be the reduced Jacobian.
Even though \(\vect{r}\) may have thousands of rows, \(\vect{J}_a\) has only
\(K\) columns.  A forward finite-difference Jacobian costs \(K+1\) residual
evaluations and is a reasonable prototype for \(K\sim20\).  Central
differences double that cost.  Colored differencing of mesh variables does not
directly help because each spline coefficient has support over several cells,
although compact B-splines preserve partial sparsity.

Better long-term options are:
\begin{enumerate}
\item analytic derivatives of the basis map combined with directional
  derivatives of the existing residual;
\item complex-step differentiation for code paths that contain no comparisons
  or nonanalytic functions of the perturbed quantity;
\item algorithmic differentiation after the equilibrium and table interfaces
  have been made compatible; or
\item matrix-free products \(\vect{J}_a\delta\coef\) in a reduced
  Gauss--Newton or trust-region method.
\end{enumerate}

If one full equilibrium residual evaluation costs \(C_R\), a finite-difference
Gauss--Newton iteration costs approximately \((K+1)C_R\), plus a small dense
least-squares solve.  This must be measured directly.  The hours-long time
integration does not determine \(C_R\); chemistry and nonlocal radiation are
expected to dominate a single call, but no runtime claim should be made before
profiling.

Warm-starting chemistry can reduce \(C_R\) during deterministic continuation.
For MCMC it is permitted only if it changes iteration count but not the final
residual.  A history-dependent approximate chemistry result is not an acceptable
speed optimization.

% =====================================================================
\section{Proposed algorithms}\label{sec:algorithms}
% =====================================================================

\subsection{Algorithm A: exact discrete reduced solve}

\begin{enumerate}
\item Select the physical branch and boundary configuration.  Decide whether
  the domain contains a transonic wind or remains subsonic at its outer edge.
\item Choose coarse spline knots and initialize \(\ln v\), \(\ln T\), and
  \(\ln\mdot\) from a cold, transonic, Wind-AE, or previously converged profile.
\item Evaluate the splines on the existing EXHALE cell centers.
\item Compute \(\rho\) from Eq.~\eqref{eq:rhoeliminate}, then obtain pressure
  from the equilibrium composition and equation of state.
\item Pack the conserved state and apply the same ghost-cell boundary routine
  used by the production code.
\item Restore a deterministic chemistry initial state and evaluate the full
  EXHALE steady residual.
\item If elemental diffusion is part of the target, either complete the current
  outer hydrodynamics--diffusion co-convergence or append the conservative
  element residual in Eq.~\eqref{eq:heweak}.  If ion-stage advection is part of
  the target, append only the independent stage equations in
  Eq.~\eqref{eq:ionweak}.
\item Scale every residual row and minimize Eq.~\eqref{eq:merit} by a
  trust-region least-squares method.
\item If progress stalls, inspect the largest physical residuals before adding
  knots.  Do not infer that the basis is too small from the scalar merit alone.
\item Increase the basis dimension by knot insertion and prolong the old spline
  into the new basis.  Continue until the solution and held-out residual are
  stable.
\item Pass the final state through the unmodified EXHALE residual diagnostic
  and, as a separate test, a short time-dependent relaxation.
\end{enumerate}

\subsection{Algorithm B: weak spectral-element solve}

\begin{enumerate}
\item Start from the converged or best available result of Algorithm A.
\item Define elements independently of the spline knots.  Place additional
  element boundaries near rapid thermal or ionization changes.
\item At Gaussian nodes, evaluate the spline state, conditional or transported
  chemistry, radiation, heating, cooling, and transport coefficients.
\item Assemble Eq.~\eqref{eq:weakresidual} for momentum and energy and, when
  active, Eqs.~\eqref{eq:heweak} and \eqref{eq:ionweak} for composition.  Mass
  is already exact by construction.
\item Add boundary and branch-appropriate sonic constraints.
\item Minimize the scaled weak merit, then repeat with larger \(Q\), more
  elements, and more spline knots.
\item Evaluate the continuous differential residual on an independent dense
  grid and the EXHALE discrete residual after projection onto its mesh.
\item Attribute differences correctly: quadrature error, basis truncation,
  continuum-versus-HLLC discretization, or a different boundary condition.
\end{enumerate}

\subsection{Algorithm C: mode search and posterior sampling}

\begin{enumerate}
\item Construct the deterministic, scaled objective and verify repeated
  evaluations first.
\item Run many inexpensive deterministic starts or an annealed population over
  the coarsest basis.
\item Cluster minima by \(\mdot\), sonic radius, temperature structure, and
  equation-wise residual, not by coefficient distance alone.
\item Refine every distinct physical mode with Algorithm A or B.
\item Only then define \(\beta\) or a physical covariance and run MCMC.
\item Confirm each sampled mode by an independent deterministic solve.
\end{enumerate}

% =====================================================================
\section{Relevant public codes and transferable methods}
\label{sec:relatedcodes}
% =====================================================================

No public code identified in this survey combines all features proposed here:
a reduced representation of a stationary astrophysical wind, the existing
EXHALE radiation and chemistry, conservative species transport, and direct
solution of the resulting coefficient residual.  There are, however, mature
public precedents for the important numerical components.  The closest
reference for the one-dimensional boundary-value algorithm is AUTO-07p.  ESTER
is the closest astrophysical precedent for a nonlinear spectral solve with
complex microphysics.  Dedalus and KADATH provide useful independent examples
of spectral equation assembly, but they are lower-priority implementation
references for EXHALE.

\begin{center}\small
\begin{tabular}{@{}p{0.17\textwidth}p{0.35\textwidth}p{0.40\textwidth}@{}}
\toprule
code & relevant capability & principal mismatch with EXHALE \\
\midrule
AUTO-07p & adaptive piecewise-polynomial collocation for ODE boundary-value
  problems; boundary and integral constraints; continuation and bifurcation
  analysis & local ODE interface; no radiation--chemistry model or existing
  finite-volume fixed point \\
ESTER & nonlinear Newton solution of rotating-star structure with radial
  Chebyshev expansions, spherical harmonics, coordinate mappings, and stellar
  microphysics & smooth two-dimensional stellar structure rather than a
  transonic irradiated wind \\
Dedalus & open spectral framework for initial-value, boundary-value, and
  eigenvalue problems; symbolic equation and boundary-condition assembly &
  adopting it as the main solver would require a separate implementation of
  EXHALE physics \\
KADATH & multidomain spectral representation and nonlinear equation systems in
  general relativity and theoretical physics & primarily elliptic equilibrium
  problems rather than conservative advection--diffusion winds \\
\bottomrule
\end{tabular}
\end{center}

\subsection{AUTO-07p: the primary numerical reference}

AUTO-07p solves ODE boundary-value problems by orthogonal collocation with
piecewise polynomials.  Its public manual specifies two to seven collocation
points in each mesh interval and adapts the mesh to equidistribute a local
discretization-error estimate.  It also treats boundary conditions, integral
constraints, parameter continuation, folds, and branch switching
\cite{Doedel2021}.  These features align directly with the reduced EXHALE
problem:
\begin{equation}
 \vect{y}(r)=
 \left(\ln v,\ln T,\operatorname{logit}f_{\rm He},N_{\rm HI},
 N_{\rm HeI},\ldots\right),
 \label{eq:autostate}
\end{equation}
while \(\ln\mdot\), a sonic radius, irradiation strength, or a diffusion
multiplier can be treated as global continuation parameters.  Nonlocal columns
become local auxiliary equations through Eq.~\eqref{eq:columnode}.

The most transferable AUTO-07p ideas are:
\begin{enumerate}
\item a local polynomial representation rather than one high-degree global
  polynomial;
\item an adaptive interval mesh driven by a defect estimate;
\item simultaneous enforcement of interior equations, physical boundary
  conditions, and global integral constraints;
\item pseudo-arclength continuation through parameter folds; and
\item explicit branch labels and restart states rather than interpreting every
  converged root as the same physical solution.
\end{enumerate}
Continuation is especially relevant to EXHALE.  A solution can first be found
with frozen heating and composition, then followed as radiation, chemistry,
conduction, or diffusion is turned on.  This is generally safer than attempting
the fully coupled root from an arbitrary initial profile.

AUTO-07p cannot accept a singular sonic equation without reformulation.  The
sonic point must be handled by a regularity condition, a desingularized
independent coordinate, or domain splitting.  Likewise, calling the complete
EXHALE finite-volume and nonlocal microphysics path from an AUTO callback would
obscure which discrete problem is being solved.  The recommended use is to
adopt its collocation, adaptivity, and continuation design in the EXHALE wrapper,
and to use the public AUTO code first on simplified Parker and frozen-source
reference problems.

\subsection{Collocation nodes are not weak-quadrature nodes}

AUTO-07p is a close algorithmic precedent, but it does not implement
Eq.~\eqref{eq:weakresidual}.  At a collocation node it enforces the strong ODE
defect,
\begin{equation}
 \vect{r}_{q}=\frac{d\vect{y}_h}{dr}(r_q)
 -\vect{f}(r_q,\vect{y}_h,\boldsymbol{\theta})=0.
 \label{eq:collocationdefect}
\end{equation}
In the weak formulation, Gaussian weights instead approximate an integral after
integration by parts.  Both methods use polynomial representations and Gaussian
nodes, but their residuals, conditioning, and conservation properties differ.
The prototype must label them separately:
\begin{enumerate}
\item \emph{strong collocation}: closest to AUTO-07p and simplest for a regular
  first-order BVP;
\item \emph{conservative weak residual}: preferred when exact element flux
  balance and discontinuous source structure dominate; and
\item \emph{existing finite-volume residual}: required when the target is the
  current EXHALE discrete fixed point.
\end{enumerate}

\subsection{ESTER: the primary astrophysical reference}

ESTER computes steady two-dimensional models of rapidly rotating stars with
spectral expansions and nonlinear iterations.  Its radial representation uses
Chebyshev polynomials, its angular representation uses spherical harmonics, and
its coordinate mapping follows the distorted stellar surface.  Newton
iterations are used when realistic microphysics makes simpler fixed-point
iterations inadequate \cite{Rieutord2016}.

ESTER demonstrates that a global astrophysical structure calculation can couple
smooth field coefficients, equation-of-state and opacity tables, radiative
transport, and large-scale flow without evolving the full transient.  The
transferable lessons are coefficient and equation scaling, continuation from
simple stellar models, mapped multidomain coordinates, and separate monitoring
of each physical equation.  Its direct spectral discretization should not be
copied blindly for an EXHALE wind: a sonic transition, outward characteristics,
and localized ionization fronts give the wind a different mathematical
structure.

\subsection{Dedalus and KADATH as independent design checks}

Dedalus is an open spectral PDE framework developed for astrophysical,
geophysical, and fluid-dynamical applications \cite{Burns2020}.  Its nonlinear
boundary-value interface rewrites symbolic equations and constraints as a root
problem and applies Newton iterations.  It is useful for rapidly testing a
simplified continuum formulation of advection, diffusion, column equations,
and boundary conditions.  It is not the recommended production dependency,
because reimplementing EXHALE chemistry, radiation, tables, and boundary logic
in a second framework would create a new physical code rather than a reduced
front end to EXHALE.

KADATH is a public C++ spectral library for nonlinear systems in theoretical
physics and numerical relativity \cite{Grandclement2010}.  Its multidomain
representations and equation-system abstraction are relevant if EXHALE later
requires separate smooth domains around a sonic point or narrow thermal front.
Its principal applications are elliptic equilibrium systems, so it supplies
software-structure and spectral-convergence precedents rather than a ready
advection--diffusion solver.

\subsection{Recommended inheritance for EXHALE}

The codes should be used as references, not combined as dependencies.  The
recommended inheritance is:
\begin{enumerate}
\item retain the EXHALE residual and microphysics ownership from Formulation A;
\item adopt AUTO-07p's piecewise-polynomial collocation, adaptive defect control,
  restart metadata, and pseudo-arclength continuation for the continuum phase;
\item adopt ESTER's scaling, progressive activation of microphysics, and mapped
  smooth-domain practices;
\item use Dedalus only for small independent equation and boundary-condition
  experiments; and
\item use KADATH as a reference if multidomain spectral geometry becomes
  necessary.
\end{enumerate}
This survey therefore changes no physical acceptance criterion.  It strengthens
the case for a deterministic adaptive collocation and continuation stage, but
the final solution must still pass the unmodified EXHALE residual or be labeled
explicitly as a continuum weak solution.

% =====================================================================
\section{Implementation plan in the EXHALE tree}\label{sec:implementation}
% =====================================================================

No production hydrodynamic path should be changed for the first prototype.
Add an isolated module and driver so that the experiment cannot alter a normal
EXHALE run.

\begin{center}\small
\begin{tabular}{@{}p{0.38\textwidth}p{0.56\textwidth}@{}}
\toprule
proposed component & responsibility \\
\midrule
\code{src/modules/time\_step/}\newline
\code{reduced\_steady\_state.f90} & spline map,
  coefficient transforms, exact mass-continuity elimination, residual scaling,
  and reduced merit \\
\code{src/utils/reduced\_steady\_driver.f90} & standalone input, initialization,
  optimization and continuation control, checkpoints, and diagnostics \\
\code{src/modules/time\_step/}\newline
\code{steady\_conservation\_residual.f90} & later implementation
  of element balances, strong collocation defects, and Gaussian weak
  quadrature; absent from Phase A \\
\code{src/modules/functions/}\newline
\code{species\_transport\_flux.f90} & common evaluation of advective,
  molecular-diffusive, and eddy-diffusive element fluxes, factored from the
  existing update so that time integration and residual assembly use identical
  physics \\
\code{tests/reduced\_steady/} & analytic winds, saved EXHALE states, basis and
  quadrature refinement gates \\
\bottomrule
\end{tabular}
\end{center}

The existing full residual evaluator currently accepts and mutates a species
state.  A reduced wrapper should expose a pure logical interface:
\begin{equation}
 (\coef,\boldsymbol{\theta})\longmapsto
 (\vect{r},\text{status},\text{diagnostics}),
 \label{eq:pureinterface}
\end{equation}
where identical inputs produce identical outputs.  The status must distinguish
invalid state, chemistry nonconvergence, boundary violation, table-range error,
and a successfully evaluated large residual.  Returning one generic penalty for
all failures creates flat regions that are difficult for both optimization and
MCMC.

Useful diagnostic output for every accepted iteration includes
\begin{itemize}
\item total and equation-wise scaled norms;
\item unscaled integrated mass, momentum, and energy imbalance;
\item element flux variation, separated advective and diffusive contributions,
  and ion-stage source-sum error when transport is active;
\item radius and equation of the largest residual;
\item \(\mdot\), sonic radius, maximum temperature, and outer Mach number;
\item chemistry iteration counts and fallback events;
\item minimum density, pressure, temperature, and species fraction; and
\item basis size, element count, and quadrature order.
\end{itemize}

Checkpoint coefficient vectors rather than only sampled profiles.  A checkpoint
must also store knot positions, coordinate mapping, residual scales, active
physics, boundary configuration, and the source revision.  Coefficients have no
meaning without that metadata.

% =====================================================================
\section{Validation program}\label{sec:validation}
% =====================================================================

\subsection{Gate 1: analytic and semi-analytic problems}

Begin with physics that has a known branch:
\begin{enumerate}
\item isothermal Parker wind in a point-mass potential;
\item hydrostatic isothermal atmosphere with \(\mdot=0\), using a formulation
  that does not divide by velocity; and
\item adiabatic wind with prescribed smooth heating.
\end{enumerate}
Measure profile error, sonic radius, conservation residual, and convergence
under basis refinement.  The Parker test is especially important because an
optimizer can find subsonic breeze and supersonic branches that have small local
residuals but do not satisfy the intended critical condition.

\subsection{Gate 2: projection of an existing EXHALE state}

Project a trusted full-grid state onto the spline basis without solving.  This
answers whether the proposed basis can represent the profile at all.  Record
profile errors and the change in the full EXHALE residual caused only by
projection.  If projection alone produces a large residual, the knots or
variables must be changed before optimizer behavior is studied.

\subsection{Gate 3: hydrodynamics with frozen microphysics}

Hold heating, cooling, and composition fixed from a reference state and solve
only the reduced hydrodynamics.  This isolates basis geometry, scaling, boundary
conditions, and sonic selection.  It is a numerical gate, not a physically
self-consistent final result.

\subsection{Gate 4: fully coupled discrete residual}

Enable equilibrium chemistry and radiation in the same path used by the current
steady solver.  Compare against a full-grid JFNK solution when one exists.
Acceptance requires agreement in the physical profiles and a small unmodified
EXHALE residual; agreement in \(\mdot\) alone is insufficient.

\subsection{Gate 5: conservative species transport}

Activate transport processes one at a time before combining them:
\begin{enumerate}
\item with diffusion disabled, verify constant elemental fraction and the
  advective element flux implied by the hydrodynamic solution;
\item with \(v=0\) and fixed temperature, recover the appropriate diffusive
  equilibrium or barometric separation profile;
\item impose the base abundance and verify zero outer diffusive flux without
  suppressing outward advective escape;
\item require \(r^2J_X\) to be constant for every conserved element and verify
  the same result from face fluxes and integrated weak residuals;
\item compare a reduced simultaneous diffusion solution with the converged
  production outer iteration, while labeling any changed energy closure;
\item in ion-advection tests, require the sum of reaction sources to conserve
  nuclei and the summed stage residual to match the element residual; and
\item compare conditional outward integration with the guarded
  \code{post\_process\_adv.f90} result only inside the latter's documented
  validity region.
\end{enumerate}
Repeat the tests in advection-dominated and diffusion-dominated limits.  A
matching abundance profile without matching total flux is not sufficient.

\subsection{Gate 6: weak formulation}

For Algorithm B, demonstrate independently:
\begin{enumerate}
\item quadrature convergence at fixed basis and elements;
\item element refinement at fixed basis;
\item basis refinement at fixed element layout;
\item small held-out strong residual;
\item stable cumulative columns and heating/cooling integrals; and
\item convergence toward Algorithm A as both discretizations are refined, or a
  documented continuum-versus-finite-volume difference.
\end{enumerate}

\subsection{Gate 7: robustness and branch search}

Use cold, transonic, Wind-AE, and perturbed starts.  A unique physical branch
should converge to the same state.  Distinct solutions must each satisfy the
same boundary conditions and residual tolerance before being labeled multiple
physical branches.  A mode caused by chemistry history, residual weighting, or
an unconverged quadrature is numerical.

\subsection{Gate 8: short dynamic check}

Map the reduced solution to the production grid and evolve it briefly with the
ordinary EXHALE time step.  The state should remain stationary within the
measured residual scale.  This check does not replace the residual test, and it
must not be extended into an hours-long regeneration merely to reproduce an
existing product.

% =====================================================================
\section{Failure modes and safeguards}\label{sec:failures}
% =====================================================================

\begin{center}\small
\begin{tabular}{@{}p{0.27\textwidth}p{0.32\textwidth}p{0.34\textwidth}@{}}
\toprule
failure & symptom & safeguard \\
\midrule
over-smooth basis & low broad merit but a missed heating or ionization layer &
  held-out residual; adaptive local knots; inspect source profiles \\
quadrature aliasing & residual small only at fitted nodes & increase \(Q\)
  without refitting; use independent nodes \\
bad residual scaling & one equation improves while another remains physically
  wrong & print equation-wise unscaled balances and largest rows \\
false Bayesian certainty & narrow coefficient posterior controlled by an
  arbitrary \(\beta\) & label it a search distribution; vary scale and basis \\
chemistry memory & repeated identical coefficients give different merit & pure
  evaluator; restore state; tighten equilibrium convergence \\
wrong branch & smooth subsonic breeze accepted as a transonic wind & impose and
  verify both sonic conditions when physically required \\
incompatible boundary conditions & nonzero residual floor localized at a
  boundary & analyze characteristics and physical BCs; do not smear the error \\
global polynomial ringing & oscillations around a narrow transition & local
  B-splines or spectral elements; knot refinement \\
coefficient degeneracy & large posterior correlations and slow MCMC & remove
  redundant basis directions; orthogonalize or whiten coefficients \\
invalid finite-volume comparison & weak solution differs from HLLC fixed point &
  report the discretization distinction; refine both formulations \\
\bottomrule
\end{tabular}
\end{center}

Transport introduces three additional failure modes.  First, separately
discretized advective and diffusive divergences can mimic an elemental source;
assemble and test one total face flux.  Second, all \(S\) ion-stage equations
plus an element equation produce a rank-deficient system; retain only \(S-1\)
independent stage constraints.  Third, setting the outer \emph{total} species
flux to zero blocks physical escape; set only the diffusive contribution to zero
when that is the intended boundary condition.  If species enthalpy transport or
frictional heating is omitted, record the approximation rather than silently
claiming conservation for a more complete multicomponent model.

A scalar \(\chi^2\) can conceal a physically wrong compromise.  In particular,
if boundary and interior constraints are incompatible, an optimizer or MCMC
sampler will distribute error according to the chosen weights.  It cannot make
the equations compatible.  Physical correctness is determined from the
individual balances, boundary characteristics, and critical-point topology,
not from a small scalar merit alone.

% =====================================================================
\section{Decision points before implementation}\label{sec:decisions}
% =====================================================================

The following choices materially change the solver and should be fixed before
production implementation:
\begin{enumerate}
\item \textbf{Target equation.} Exact current finite-volume fixed point
  (Formulation A), continuum weak solution (Formulation B), or both in sequence.
\item \textbf{Physical branch.} Outward transonic wind, subsonic L1-truncated
  flow, or hydrostatic atmosphere.  One parameterization should not silently
  force all three.
\item \textbf{Chemistry closure.} Local equilibrium eliminated at each trial,
  or advected ion fractions represented as additional fields.
\item \textbf{Transport closure.} Production operator splitting with outer
  co-convergence, or a simultaneous conservative residual for elemental
  diffusion and ion-stage advection.  Also state whether diffusive enthalpy and
  frictional heating belong to the target model.
\item \textbf{Boundary treatment.} Hard constraints built into the basis versus
  soft residuals with stated scales.
\item \textbf{Role of probability.} Annealed global search, numerical
  discretization uncertainty, uncertain physical parameters, or observational
  inference.  These interpretations require different covariance models.
\end{enumerate}

For the current EXHALE structure, the recommended first decision is narrow:
target the discrete residual, solve an outward wind with equilibrium chemistry,
hard-enforce base temperature and steady mass continuity, retain the production
boundary routine, and use deterministic trust-region least squares.  This
prototype directly tests the central hypothesis: that a smooth 20--30 dimensional
state manifold can reach a fixed point that the full \(3N\)-variable solver has
difficulty approaching.  Elemental diffusion should then be added first by
reproducing the production outer co-convergence, followed by a simultaneous
conservative transport residual.  Self-consistent ion-stage advection is a
later extension because it changes the chemistry closure and coefficient count.
Once this discrete prototype is stable, introduce AUTO-style adaptive
collocation and pseudo-arclength continuation without changing the production
microphysics implementation.

% =====================================================================
\section{Acceptance criteria for a successful prototype}\label{sec:acceptance}
% =====================================================================

The prototype is successful only if all of the following are demonstrated:
\begin{enumerate}
\item repeated coefficient evaluations are deterministic to the requested
  tolerance;
\item the reduced solution satisfies mass continuity by construction and both
  momentum and energy residuals under the unmodified production diagnostic;
\item when species transport is active, total element fluxes satisfy their
  conservative residuals, transport boundary conditions are met, and ion-stage
  source sums conserve nuclei;
\item results are stable under basis refinement and initialization;
\item the inferred sonic structure and outer characteristic state are physically
  consistent with the selected domain;
\item the final profile remains stationary in a short production-code evolution;
\item any difference from a reference solution is separated into physics,
  boundary, basis, quadrature, and discretization contributions; and
\item measured wall-clock cost, residual-call count, and chemistry failures are
  reported without extrapolating from the hours-long marching runtime.
\end{enumerate}

MCMC is an optional success criterion.  It becomes useful after the deterministic
solver works, either because several valid branches exist or because uncertain
physical inputs and observations define a genuine posterior.  Failure of a
random-walk chain to mix is not evidence that the reduced steady formulation is
wrong; failure of the independently evaluated physical residual is.

% =====================================================================
\section{Summary}\label{sec:summary}
% =====================================================================

The proposed method replaces the full radial state by smooth spline expansions
of velocity and temperature plus one mass-loss parameter.  Density is eliminated
exactly through steady mass continuity, while the existing EXHALE equilibrium
chemistry determines composition and pressure conditionally.  The first solver
should minimize the current full-grid finite-volume residual, because it gives
an unambiguous comparison with the production fixed point.  A later weak
spectral-element solver can integrate conservative momentum and energy balances
with Gaussian quadrature and may require far fewer radial evaluations.

Advection and diffusion fit the same framework.  Bulk hydrodynamic advection,
viscosity, and conduction are already part of the current discrete residual.
Elemental diffusion requires a small set of additional abundance coefficients
and a residual for the total advective--diffusive element flux.  Ion-stage
advection requires either independent softmax spline fields or a conditional
upstream integration.  The current operator-split diffusion and post-processed
ion advection must not be described as if they were already simultaneous JFNK
residual blocks.

The production workflow is more robust to sharp structure, has an unambiguous
conservative finite-volume fixed point, and provides a fallback when Newton
fails.  Its main costs are the long transient, the high-dimensional JFNK state,
and split treatment of composition transport.  The reduced workflow removes
many redundant search directions and supports direct continuation, but adds
basis error, explicit sonic-branch responsibility, and a possible distinction
from the production discretization.  The recommended architecture therefore
uses the reduced solve to enter the full-grid JFNK basin and retains ordinary
EXHALE as the final residual diagnostic and short stability test.

Among the surveyed public codes, AUTO-07p is the closest numerical reference
for the one-dimensional reduced boundary-value problem because it combines
piecewise-polynomial collocation, adaptive meshes, boundary and integral
constraints, and continuation.  ESTER supplies the closest astrophysical
precedent for a nonlinear spectral solve coupled to realistic microphysics.
Dedalus is useful for independent continuum prototypes, while KADATH becomes
relevant only if a multidomain spectral representation is required.  These
precedents support the deterministic solver architecture but do not replace the
EXHALE conservation and fixed-point validation gates.

The scalar residual may be written in a likelihood form, but nonlinear least
squares is the appropriate first solver.  Annealed MCMC or sequential Monte
Carlo can search for multiple branches; posterior MCMC is meaningful when a
physical error model or observational likelihood is present.  In every case,
equation-wise conservation, boundary consistency, deterministic chemistry,
basis refinement, quadrature refinement, and held-out residuals remain the
acceptance criteria.

\bibliographystyle{plain}
\bibliography{reduced_basis_mcmc_steady_solver}

\end{document}
